\section{System Design and Architecture}
\label{sec:design}

\subsection{Design Goals}

\software{} was designed to satisfy six key requirements distilled from the
experience of the authors and feedback from early users in the research
software community:

\begin{description}
  \item[RG1: Local-first project model.] Manuscript content must reside on the
  researcher's local filesystem as ordinary directories and files, compatible
  with version control (git), external tools (grep, diff), and manual
  inspection. The system should augment---not replace---existing file-based
  workflows.

  \item[RG2: Permission-aware AI interaction.] AI assistance must be separated
  into distinguishable modes with clear boundaries: read-only discussion,
  reviewable edit proposals with inline diff visualization, and controlled tool
  execution. Users must be able to inspect and approve every file change before
  it is applied.

  \item[RG3: Verifiable artifact production.] The system must integrate
  citation verification and \LaTeX{} compilation so that generated manuscripts
  can be checked for reference accuracy and build correctness, not just text
  fluency.

  \item[RG4: Auditable workflow pipelines.] Multi-stage writing workflows must
  be executable with human checkpoints at decision points, producing a record
  of what AI stages were executed, what human feedback was provided, and what
  artifacts were produced at each stage.

  \item[RG5: Interoperability.] The system must expose core capabilities
  through standard protocols so that external tools and MCP-compatible clients
  can integrate with the platform.

  \item[RG6: Privacy and configuration safety.] API keys and sensitive
  configuration must be stored in local environment files, never exposed to the
  browser, and masked in API responses.
\end{description}

Figure~\ref{fig:workspace-screenshot} highlights the implemented workspace
features that distinguish \software{} from a general-purpose editor or chatbot.
The left pane keeps the filesystem-backed project visible as the source of
truth, the center pane couples \LaTeX{} source editing with compiled PDF
feedback, and the right pane exposes AI assistance and typed pipeline controls
inside the same local project context.

\begin{figure*}[!b]
  \centering
  \includegraphics[width=\textwidth]{fig-workspace-screenshot.pdf}
  \caption{Paper Agent workspace during manuscript editing, emphasizing the
  system's main differentiators: a filesystem-backed project tree, synchronized
  source/PDF editing, local manuscript artifacts, and typed Pipeline 2.0 stages
  for reviewable AI-assisted workflows. The screenshot shows these capabilities
  co-located rather than separated across a standalone chatbot, editor, and build
  tool.}
  \label{fig:workspace-screenshot}
\end{figure*}

\subsection{Architecture Overview}

\software{} follows a client-server architecture with a filesystem-backed
project model. Figure~\ref{fig:architecture} provides a high-level overview of
the system components and their interactions.

\begin{figure}[htbp]
  \centering
  \includegraphics[width=\textwidth]{fig-architecture.pdf}
  \caption{System architecture overview. The system follows a client-server
  model with four layers: presentation (React/Vite), application
  (Fastify/Node.js), service (AI providers, citation APIs, \LaTeX{} engines),
  and storage (local filesystem). External MCP clients can invoke registered
  tools through the JSON-RPC/SSE transport.}
  \label{fig:architecture}
\end{figure}

The system comprises four main layers:

\textbf{Presentation layer} (React 18 / Vite / TypeScript). A three-panel
browser interface provides a project file tree, a multi-mode editor (source,
split, rendered) with CodeMirror 6, \LaTeX{} preview via KaTeX, PDF viewing
via the browser's native PDF viewer, and a right panel with tabs for AI Chat, Skills (40+ writing
plugins), Review, Anti-AI detection, and Pipeline 2.0 controls.

\textbf{Application layer} (Fastify / Node.js). REST API endpoints serve
project management, file operations, AI conversation streaming (via SSE),
citation verification, \LaTeX{} compilation, pipeline orchestration, and MCP
tool exposure. The application layer also manages tmux terminal sessions for
integrated shell access.

\textbf{Service layer.} Backend services encapsulate AI provider integration
(OpenAI/Anthropic-compatible), citation verification (CrossRef REST API,
Semantic Scholar API, OpenAlex API), \LaTeX{} compilation engines (pdflatex,
xelatex, lualatex, latexmk, tectonic), tmux terminal management, and pipeline
stage execution.

\textbf{Storage layer.} Manuscript projects are stored as directories under a
configurable \texttt{papers/} path. Each project contains \LaTeX{} source
files, a \texttt{project.json} metadata file, and optional subdirectories
(\texttt{fig/}, \texttt{docs/}, \texttt{code/}). Configuration lives in a
repository-root \texttt{.env} file, loaded at startup and writable through the
API with masked-key responses.

\subsection{Key Design Decisions}

\subsubsection{Filesystem as the Source of Truth}

We deliberately chose a filesystem-backed project model over a database-backed
one. This decision means that projects are ordinary directories---they can be
versioned with git, edited with external tools (VS Code, TeXstudio, vim),
backed up with standard file operations, and shared without depending on
\software{}. The tradeoff is that the backend must handle concurrent
filesystem access and cannot rely on database transactions for consistency. We
mitigate this by using atomic file writes (write to temporary file, then
rename), validating file state before operations, and operating under the
assumption that manuscript editing is inherently sequential at the project
level.

\subsubsection{Separation of AI Permission Modes}

Rather than a binary ``AI on/off'' toggle, we designed three modes with
escalating capabilities, illustrated in Figure~\ref{fig:ai-modes}.

\begin{figure}[htbp]
  \centering
  \includegraphics[width=\textwidth]{fig-ai-modes.pdf}
  \caption{Permission-aware AI interaction modes. Chat mode provides read-only
  discussion with no file modifications. Agent mode generates unified diff
  proposals that require explicit user approval before application. Tools mode
  enables multi-step operations constrained to the project code directory with
  full audit logging.}
  \label{fig:ai-modes}
\end{figure}

\begin{itemize}
  \item \textbf{Chat mode:} Read-only discussion. The AI receives the
  manuscript context (chapter content, references) but no file-writing or
  code-execution tools. Output is plain text responses, and no state change
  occurs outside the conversation.
  \item \textbf{Agent mode:} Can propose edits as unified diffs with inline
  color-coded visualization. The frontend renders add/del lines side by side
  with Accept/Reject buttons. No change is applied without user approval.
  \item \textbf{Tools mode:} Full multi-step tool execution, constrained to
  the project's \texttt{code/} directory. All tool calls and results are
  displayed in the conversation for audit. This mode is designed for
  operations such as running experiments, generating figures, or formatting
  bibliographies.
\end{itemize}

This design emerged from early development experience showing that researchers
wanted to discuss ideas freely (Chat) before committing to specific changes.
Research on human-in-the-loop machine learning~\cite{hitl_survey} has
identified lack of control over AI-generated changes as a top concern across
AI-assisted tasks, confirming our design direction.

\subsubsection{Pipeline as Typed Stages with Human Gates}

Early versions of \software{} used a linear pipeline model where stages
executed sequentially. We found that different stages require fundamentally
different execution models: AI stages need LLM context management and
streaming, compute stages need shell timeout and error handling, citation
stages need API rate limiting with exponential backoff, and compile stages
need error parsing with log analysis. Figure~\ref{fig:pipeline} illustrates
the stage orchestration.

\begin{figure}[htbp]
  \centering
  \includegraphics[width=\textwidth]{fig-pipeline.pdf}
  \caption{Pipeline 2.0 typed stage system. Five stage types (AI, Compute,
  Human, Citation, Compile) are orchestrated with human checkpoints as
  first-class gates. Failed compile stages route back to AI for revision, while
  successful stages produce verified PDF artifacts.}
  \label{fig:pipeline}
\end{figure}

Pipeline 2.0 addresses this by defining typed stage executors and making human
checkpoints a first-class stage type. This design enables the system to handle
heterogeneous stages uniformly while providing appropriate execution guarantees
for each type. Five preset pipeline templates are provided:
\begin{itemize}
  \item \textbf{Writing Flow:} Outline $\to$ Draft $\to$ Polish $\to$ Review
  \item \textbf{Paper Pipeline:} Polish $\to$ Review $\to$ Revise $\to$
  Citation Check $\to$ Compile
  \item \textbf{Quick Review:} Review $\to$ Revise
  \item \textbf{Citation Pipeline:} Verify $\to$ Deduplicate $\to$ Discover
  \item \textbf{Executable Paper:} Run Code $\to$ Generate Figures $\to$ Compile
\end{itemize}

\subsubsection{MCP as the Interoperability Boundary}

Rather than designing a custom API for external integration, we adopted the
Model Context Protocol~\cite{modelcontextprotocol2024}---an open standard from
Anthropic for tool-exposing servers. The MCP server is implemented as a
JSON-RPC 2.0~\cite{jsonrpc} endpoint over SSE transport, registering tools
for paper search, citation verification, \LaTeX{} compilation, file reading,
and AI operations. This decision means that \software{} tools are immediately
usable from any MCP-compatible client (Claude Desktop, Cursor, GitHub Copilot)
without additional integration work.
